% -*- coding: utf-8 -*-
\documentclass[aspectratio=169,10pt,UTF8]{ctexbeamer}

% -------------------------------------------------------------
% 主题与宏包配置 (Academic Clean Metropolis / Madrid Style)
% -------------------------------------------------------------
\usetheme{Madrid}
\usecolortheme{default}

\usepackage{amsmath,amssymb,amsfonts}
\usepackage{booktabs}
\usepackage{tikz}
\usepackage{pgfplots}
\pgfplotsset{compat=1.18}
\usepackage{hyperref}

% 配色方案定义 (现代学术深蓝与石板灰)
\definecolor{AcademiaBlue}{RGB}{24, 76, 120}
\definecolor{AccentBlue}{RGB}{41, 128, 185}
\definecolor{DarkSlate}{RGB}{44, 62, 80}
\definecolor{LightBg}{RGB}{248, 250, 252}

\setbeamercolor{structure}{fg=AcademiaBlue}
\setbeamercolor{title}{fg=white,bg=AcademiaBlue}
\setbeamercolor{frametitle}{fg=white,bg=AcademiaBlue}
\setbeamercolor{block title}{fg=white,bg=AcademiaBlue}
\setbeamercolor{block body}{bg=LightBg,fg=DarkSlate}

% -------------------------------------------------------------
% 报告元数据 (Metadata)
% -------------------------------------------------------------
\title[博士文献汇报：前沿综述与开题切入点]{面向[核心课题方向]的深度学术调研与前沿综述}
\subtitle{组会文献精读分享 / 博士学位论文开题文献调研}
\author[汇报人]{汇报人：博士研究生 \quad 指导教师：导师姓名}
\institute[实验室名称]{所属学院 / 重点实验室 / 课题组}
\date{\today}

\begin{document}

% 幻灯片 1: 封面
\begin{frame}
  \titlepage
\end{frame}

% 幻灯片 2: 大纲
\begin{frame}{汇报大纲 (Agenda)}
  \tableofcontents
\end{frame}

% -------------------------------------------------------------
\section{研究背景与核心矛盾}
% -------------------------------------------------------------
\begin{frame}{1. 研究背景与核心科学问题 (Motivation)}
  \begin{block}{核心技术矛盾定义}
    现有主流模型在处理 \textbf{[目标任务]} 时，由于面临 \textbf{[理论/算力瓶颈]}，导致发生不可逆的 \textbf{[失效模式]}。
  \end{block}
  \vspace{0.3cm}
  \begin{itemize}
    \item \textbf{理论价值}：关乎泛化误差下界 $\mathcal{O}(1/\sqrt{n})$ 的紧致性证明与表征解耦。
    \item \textbf{工程落地}：直接制约了高吞吐、低显存边缘端设备的实时部署能力。
    \item \textbf{传统方法瓶颈}：在分布偏移 (OOD) 下性能衰减达 30\% 以上。
  \end{itemize}
  \note{汇报提示：用1-2分钟阐述现实背景与科学瓶颈，强调本研究方向为什么非做不可。}
\end{frame}

% -------------------------------------------------------------
\section{领域技术演进脉络}
% -------------------------------------------------------------
\begin{frame}{2. 领域技术发展脉络 (Timeline \& Evolution)}
  \begin{columns}[T]
    \begin{column}{0.32\textwidth}
      \begin{alertblock}{第一阶段 (2015-2018)}
        \textbf{浅层与传统规则}
        \begin{itemize}
          \item 依赖人工设计先验
          \item 表达容量受限
          \item 泛化边界窄
        \end{itemize}
      \end{alertblock}
    \end{column}
    \begin{column}{0.32\textwidth}
      \begin{block}{第二阶段 (2019-2022)}
        \textbf{经典深度神经网络}
        \begin{itemize}
          \item CNN / 基础 Transformer
          \item 经验精度显著攀升
          \item 计算开销与显存激增
        \end{itemize}
      \end{block}
    \end{column}
    \begin{column}{0.32\textwidth}
      \begin{exampleblock}{第三阶段 (2023-至今)}
        \textbf{前沿自监督与稀疏范式}
        \begin{itemize}
          \item 扩散模型 / SSM 架构
          \item 追求极高计算能效比
          \item 理论泛化可解释性
        \end{itemize}
      \end{exampleblock}
    \end{column}
  \end{columns}
\end{frame}

% -------------------------------------------------------------
\section{核心代表文献深度剖析}
% -------------------------------------------------------------
\begin{frame}{3. 重点文献精读：[Paper Title A]}
  \small
  \textbf{文献来源}：Author et al. (\textit{NeurIPS 2024}, Stanford University)
  \begin{columns}[T]
    \begin{column}{0.5\textwidth}
      \begin{block}{核心方法设计}
        \begin{itemize}
          \item 提出稀疏路由机制，打破注意力二次方开销 $\mathcal{O}(N^2)$。
          \item 优化目标：$\mathcal{L} = \mathcal{L}_{task} + \lambda \mathcal{L}_{sparse}$。
          \item 动态聚类 Token 减少长程冗余信息传递。
        \end{itemize}
      \end{block}
    \end{column}
    \begin{column}{0.5\textwidth}
      \begin{block}{量化实证与局限}
        \begin{itemize}
          \item \textbf{实验结果}：ImageNet 256x256 上 FID 降至 2.15 (提升 0.42)。
          \item \textbf{吞吐量}：提高 2.3x，显存下降 45\%。
          \item \textbf{学术局限}：聚类超参对扰动敏感；小目标特征丢失；\textit{文献未说明预热策略}。
        \end{itemize}
      \end{block}
    \end{column}
  \end{columns}
  \note{汇报提示：重点剖析作者为何要引入该稀疏项，并指出其消融实验中的未解释细节。}
\end{frame}

% -------------------------------------------------------------
\section{横向对比与技术矩阵}
% -------------------------------------------------------------
\begin{frame}{4. 代表性方法横向对比矩阵 (三线表规范)}
  \centering
  \scriptsize
  \begin{tabular}{lcccccc}
    \toprule
    \textbf{方法代表} & \textbf{发表年份} & \textbf{核心技术范式} & \textbf{时间复杂度} & \textbf{标注监督量} & \textbf{OOD泛化} & \textbf{本课题定位} \\
    \midrule
    Method A (Zhang) & 2024 (NeurIPS) & 稀疏路由注意力 & $\mathcal{O}(N \log N)$ & 纯自监督 & 中等 & \textbf{骨干改进基线} \\
    Method B (Wang)  & 2023 (ICML)    & 正交特征解耦   & $\mathcal{O}(N^2)$      & 弱监督   & 极强 & 对照 SOTA \\
    Method C (Chen)  & 2024 (CVPR)    & 扩散隐空间建模 & $\mathcal{O}(T \cdot N)$& 多模态对齐 & 强   & 实验竞品基准 \\
    \bottomrule
  \end{tabular}
  \vspace{0.4cm}
  \begin{itemize}
    \item \textbf{已解决瓶颈}：标准分布下的经验平均拟合精度趋于饱和。
    \item \textbf{未解决瓶颈}：极端噪声与长尾样本下的表征坍塌；紧致理论泛化界缺失。
  \end{itemize}
\end{frame}

% -------------------------------------------------------------
\section{未来挑战与博士开题切入点}
% -------------------------------------------------------------
\begin{frame}{5. 博士学位论文选题思考与切入路线}
  \begin{block}{拟定博士研究方向：[具体选题名称]}
    \begin{itemize}
      \item \textbf{科学猜想 (Hypothesis)}：在稀疏路由机制中注入因果不变性先验，可以在维持 $\mathcal{O}(N)$ 计算复杂度的同时消除虚假关联。
      \item \textbf{实施方案 (Roadmap)}：
        \begin{enumerate}
          \item 搭建开源基线框架并严格复现 Method A 实验。
          \item 设计新型因果解耦模块，推导反事实约束优化目标。
          \item 开展完整消融实验与压力测试，完成理论上界证明。
        \end{enumerate}
      \item \textbf{预期成果}：冲击 CCF-A 类顶级会议 / 中科院一区顶刊。
    \end{itemize}
  \end{block}
  \note{汇报提示：向导师展现严密的科研演进逻辑，从前人缺陷自然推导出你的立项依据。}
\end{frame}

% 结束页
\begin{frame}[plain]
  \begin{center}
    \Huge \textcolor{AcademiaBlue}{\textbf{Q \& A}} \\
    \vspace{0.5cm}
    \normalsize 请各位老师与同组同学批评指正！
  \end{center}
\end{frame}

\end{document}
